\documentclass[11pt]{article}
\usepackage[a4paper,margin=18mm]{geometry}
\usepackage{fontspec}
\setmainfont{Latin Modern Roman}
\usepackage{amsmath,amssymb,amsthm,amscd,graphicx,color}
\usepackage{xurl}
\usepackage[unicode,hidelinks]{hyperref}
\allowdisplaybreaks
\setlength\emergencystretch{3em}
\usepackage{amscd,array,enumerate}

\newcommand{\onto}{\to\!\!\!\!\!\to}

\numberwithin{equation}{section}
\newtheorem{Theorem}{Theorem}[section]
\newtheorem{Conjecture}[Theorem]{Conjecture}
\newtheorem{Corollary}[Theorem]{Corollary}
\newtheorem{Claim}[Theorem]{Claim}
\newtheorem{Maintheorem}[Theorem]{Main Theorem}
\newtheorem{Lemma}[Theorem]{Lemma}
\newtheorem{Proposition}[Theorem]{Proposition}
{ \theoremstyle{definition}
\newtheorem{Definition}[Theorem]{Definition}
\newtheorem{Example}[Theorem]{Example}
\newtheorem{Remark}[Theorem]{Remark}
}

\makeatletter\newlabel{subsec:shuffle_gen_feigin}{{6.1}{}{Original source locator}{}{}}\makeatother
\begin{document}
\begin{center}\Large Quantum flag-counting shuffle character is a twisted bialgebra homomorphism for finite-field acyclic valued quivers\end{center}
\noindent\textbf{Stable ID:} 2006\par
\noindent\textbf{Authors:} Dylan Rupel\par
\noindent\textbf{Original work:} The Feigin Tetrahedron\par
\noindent\textbf{Version:} Official SIGMA 11 (2015) article 024, DOI 10.3842/SIGMA.2015.024; journal PDF and native TeX acquired 2026-10-01, exact bytes pinned by SHA256\par
\noindent\textbf{Selected task scope:} Full Theorem7.1 under the standing finite-field and acyclic valued-quiver hypotheses: formula(7.2) defines a grading-preserving homomorphism from the graded dual Hall-Ringel bialgebra into the paper quantum shuffle bialgebra, with the displayed grading-twisted tensor products. No generic injectivity or conjectural cluster/canonical-basis assertion.\par
\noindent\textbf{Difficulty:} H2: The proof turns weighted counts of composition flags into a map between the dual Hall-Ringel and quantum shuffle bialgebras. Multiplicativity requires matching Green-theorem flag counts with all shuffle inversions and noninversions, while compatibility with coproduct requires a separate concatenation/flag identity and careful rescaling of the dual Hall basis.\par
\noindent\textbf{Editorial boundary note (not author text):} The complete original quantum-shuffle-character Theorem 7.1 and all author flag-count proof are supplied, with finite-field valued-quiver, adjacency/Euler–Ringel forms, exact half-power shuffle conventions, dual-Hall rescaling and grading-twisted tensor products. Full shuffle associativity/coproduct compatibility, all multiplicative flag/inversion-exponent identities and the separate concatenation/coproduct flag identity are retained. The original tetrahedron remains an original diagram asset; mathematical proof text and formulas are editable native TeX. Ringel/Green compatibility is cited established background, without independent reproof. Source tensor/summation symbols are preserved with the existing notice. No injectivity, Feigin-factorization companion, KLR or conjectural cluster-monomial conclusion is selected.\par
\noindent\textbf{Source:} \url{https://sigma-journal.com/2015/024/}\quad DOI: \url{https://doi.org/10.3842/SIGMA.2015.024}\par
\noindent\textbf{License and attribution:} Copyright retained by the named authors. Published by SIGMA under CC BY 4.0: \url{https://creativecommons.org/licenses/by/4.0/}. Source: \url{https://sigma-journal.com/about.html}. Adaptation consists of standalone excerpt formatting, cover and numbering restoration; original author language and mathematical text are retained.\par
\noindent\textbf{Editorial symbol notice (not author text):} In the flag-count proof some displayed tensor factors or summation subscripts use V where the surrounding coefficient comparison fixes W or a pair(U,W); the author subsequent explicit exponent and concatenation identities specify the operands. Source bytes remain unchanged.\par
\noindent\textbf{External original locator subsec:shuffle\_gen\_feigin:} Original Section 6.1 treats the unselected shuffle-to-Feigin factorization, used only in the introductory tetrahedron overview.\par
\medskip\hrule\medskip\noindent\textbf{Author text begins below.}\par
\setcounter{section}{1}
\setcounter{subsection}{0}
\setcounter{subsubsection}{0}
\setcounter{Theorem}{0}
\setcounter{equation}{0}
% Original source lines 184--281
\section{Combinatorial conventions and notations}\label{sec:combinatorics}

Fix an indeterminate~$q$.
We def\/ine the~$q$-numbers,~$q$-factorials, and~$q$-binomials~by
\begin{gather*}
(n)=(n)_q=1+q+\cdots+q^{n-1},
\\
(n)!=(n)\cdot(n-1)\cdots(2)\cdot(1),
\qquad
{n\choose k}_q=\frac{(n)!}{(k)!\cdot(n-k)!}.
\end{gather*}
In def\/ining our generalized Feigin homomorphism $\overline\Psi_\mathbf{i}$ we will make particular use of the
bar-invariant versions of the quantum numbers, for convenience we introduce $v=\sqrt{q}$:
\begin{gather*}
[n]=[n]_q=v^{-n+1}+v^{-n+3}+\cdots+v^{n-3}+v^{n-1},
\\
[n]!=[n]\cdot[n-1]\cdots[2]\cdot[1],
\qquad
{n\brack k}_q=\frac{[n]!}{[k]!\cdot[n-k]!}.
\end{gather*}
The~$q$-binomials satisfy several useful analogues of classical binomial identities.
We collect these and sketch their proofs in the following.
\begin{Lemma}\label{lem:binomial_identities}\qquad
\begin{enumerate}\itemsep=0pt
\item[$1.$] The~$q$-binomial coefficients satisfy the following identities:
\begin{enumerate}[\upshape (a)]\itemsep=0pt
\item Pascal identities: ${n\choose k}_q={n-1\choose k-1}_q+q^k{n-1\choose k}_q=q^{n-k}{n-1\choose k-1}_q+{n-1\choose k}_q$;
\item row identity: for $n>0$, $\sum\limits_{k=0}^n (-1)^k q^{-nk+\frac{1}{2} k(k+1)}{n\choose k}_q=0$;
\item subspace identity: ${m+n\choose k}_q=\sum\limits_{r+s=k}q^{r(n-s)}{m\choose r}_q{n\choose s}_q$.
\end{enumerate}
\item[$2.$] The bar-invariant~$q$-binomial coefficients satisfy the following identities:
\begin{enumerate}[\upshape (a)]\itemsep=0pt
\item Pascal identities: ${n\brack k}_q=v^{n-k}{n-1\brack k-1}_q+v^{-k}{n-1\brack k}_q=v^{k-n}{n-1\brack k-1}_q+v^k{n-1\brack k}_q$;
\item row identity: for $n>0$, $\sum\limits_{k=0}^n (-1)^k v^{k-nk} {n\brack k}_q=0$;
\item subspace identity: ${m+n\brack k}_q=\sum\limits_{r+s=k}v^{r(n-s)-s(m-r)}{m\brack r}_q{n\brack s}_q$.
\end{enumerate}
\end{enumerate}
\end{Lemma}

\begin{proof}
We prove the identities in (1).
The Pascal identities in (a) follow from the def\/inition and the equalities
$(n)_q=(k)_q+q^k(n-k)_q=q^{n-k}(k)_q+(n-k)_q$.
For (b) we simply apply the second Pascal identity to get a~telescopic summation:
\begin{gather*}
\sum\limits_{k=0}^n (-1)^k q^{-nk+\frac{1}{2} k(k+1)}{n\choose k}_q
\\
\qquad
=\sum\limits_{k=1}^n (-1)^k
q^{-n(k-1)+\frac{1}{2}(k-1)k}{n-1\choose k-1}_q+\sum\limits_{k=0}^{n-1} (-1)^k q^{-nk+\frac{1}{2} k(k+1)}{n-1\choose
k}_q=0.
\end{gather*}
To see (c) we consider a~f\/inite f\/ield $\mathbb{F}$ with~$q$ elements.
Write $\operatorname{Gr}_k(\mathbb{F}^n)$ for the set of~$k$-dimensional subspaces of $\mathbb{F}^n$.
The number of points in $\operatorname{Gr}_k(\mathbb{F}^n)$ is given by ${n\choose k}_q$.
Fix a~decomposition $\mathbb{F}^{m+n}=\mathbb{F}^m\oplus\mathbb{F}^n$ so that any subspace of $\mathbb{F}^{m+n}$ can be
decomposed as a~direct sum of a~subspace of $\mathbb{F}^m$ and a~subspace of $\mathbb{F}^n$.
This gives rise to a~surjective map $\operatorname{Gr}_k(\mathbb{F}^{m+n})\onto \bigsqcup\limits_{r+s=k}\operatorname{Gr}_r(\mathbb{F}^m)\times\operatorname{Gr}_s(\mathbb{F}^n)$ with f\/iber over any point
of $\operatorname{Gr}_r(\mathbb{F}^m)\times\operatorname{Gr}_s(\mathbb{F}^n)$ an af\/f\/ine space of dimension $r(n-s)$.
Now counting points completes the proof when~$q$ is a~power of a~prime.
But this is an equality of polynomials for inf\/initely many values and thus the polynomials must be equal for every~$q$.

The identities in (2) easily follow from those in (1) using the relation ${n\brack k}_q=v^{-k(n-k)}{n\choose k}_q$.
\end{proof}

We choose a~notation for symmetric groups which is most convenient for describing the structure constants of the
multiplication in the quantum shuf\/f\/le algebra.
For each positive integer~$t$ let $\Sigma_t$ denote the symmetric group on~$t$ letters, thought of as the set of all
orderings $\sigma=(\sigma_1,\ldots,\sigma_t)$ of the set $(1,\ldots,t)$.
We write $\sigma\cdot\tau=(\sigma_{\tau_1},\ldots,\sigma_{\tau_t})$ for the multiplication in $\Sigma_t$.
Let $\sigma^{-1}_k$ denote the position of~$k$ in~$\sigma$, i.e the number~$i$ so that $\sigma_i=k$.

For positive integers~$r$ and~$s$ def\/ine the \emph{shuffle subgroup} $\Sigma_{r,s}\subset\Sigma_{r+s}$ consisting of all
shuf\/f\/les of the sequences $(1,\ldots,r)$ and $(r+1,\ldots,r+s)$, that is
\begin{gather*}
\Sigma_{r,s}=\big\{\sigma\in\Sigma_{r+s}: \sigma^{-1}_1<\cdots<\sigma^{-1}_r, \sigma^{-1}_{r+1}<\cdots<\sigma^{-1}_{r+s}\big\}.
\end{gather*}
By convention we take $\Sigma_{0,0}=\Sigma_0=\{1\}$ to be the trivial group.

\section{Quantum groups and Feigin homomorphisms}\label{sec:feigin}

In this section we introduce the main objects underlying the present investigations: the quantum groups associated to
a~symmetrizable Cartan matrix and the Feigin homomorphisms mapping them to quantum polynomial algebras.

Fix an index set~$I$.
Let $A=(a_{ij})$ be an $I\times I$ Cartan matrix with symmetrizing matrix $D=\operatorname{diag}(\mathbf{d})$, where
$\mathbf{d}=(d_i: i\in I)$ is an~$I$-tuple of positive integers.
More explicitly, the entries of~$A$ have the following properties:
\begin{itemize}\itemsep=0pt
\item $a_{ii}=2$ for all~$i$;
\item $a_{ij}\le0$ for $i\ne j$;
\item $d_ia_{ij}=d_ja_{ji}$ for all $i,j\in I$.
\end{itemize}
Let $\mathcal{Q}$ denote the root lattice of a~Kac--Moody root system~$\Phi$ associated to~$A$.
Denoting by $\{\alpha_i\}_{i\in I}$ the simple roots of~$\Phi$, $\mathcal{Q}$ is the free abelian %Abelian ?
group with basis $\{\alpha_i\}$.
The pair $(A,\mathbf{d})$ determines a~symmetric bilinear form $(\cdot,\cdot):\mathcal{Q}\times\mathcal{Q}\to\mathbb{Z}$
given on generators of $\mathcal{Q}$ by $(\alpha_i,\alpha_j)=d_ia_{ij}$ for $i,j\in I$.
\setcounter{section}{4}
\setcounter{subsection}{0}
\setcounter{subsubsection}{0}
\setcounter{Theorem}{0}
\setcounter{equation}{0}
% Original source lines 373--390
Let~$Q$ be an acyclic quiver with vertex set~$I$ and recall the symmetrizing matrix $D=\operatorname{diag}(\mathbf{d})$,
where $\mathbf{d}=(d_i: i\in I)$.
Write $n_{ij}=n_{ji}$ for the number of arrows connecting vertices $i,j\in I$ in~$Q$ and note that all such arrows point in the same direction.
We assume further that there are only f\/initely many arrows whose source or target is the vertex $i\in I$.
From the pair $(Q,\mathbf{d})$, called a~\emph{valued quiver}, we def\/ine an adjacency matrix $B=B_Q=(b_{ij})$ by the rule
\begin{gather*}
b_{ij}=
\begin{cases}
n_{ij}d_j/\gcd(d_i,d_j) & \text{if $i\to j$ in~$Q$,}
\\
-n_{ij}d_j/\gcd(d_i,d_j) & \text{if $j\to i$ in~$Q$,}
\\
0 & \text{if $i=j$.}
\end{cases}
\end{gather*}
Notice that the matrix~$B$ is skew-symmetrizable with skew-symmetrizing matrix~$D$, i.e.~$d_ib_{ij}=-d_jb_{ji}$ for all $i,j\in I$.
To connect with the previous section we will assume that the matrix~$B$ is related to the Cartan matrix~$A$~by
$a_{ij}=-|b_{ij}|$ for $i\ne j$.
\setcounter{section}{4}
\setcounter{subsection}{0}
\setcounter{subsubsection}{0}
\setcounter{Theorem}{2}
\setcounter{equation}{0}
% Original source lines 468--510
\section{Representations of valued quivers\\ and their Hall--Ringel
algebras}\label{sec:valued quivers}

Fix a~f\/inite f\/ield $\mathbb{F}$ and an algebraic closure $\bar{\mathbb{F}}$.
For each positive integer~$d$ write $\mathbb{F}_d$ for the degree~$d$ extension of $\mathbb{F}$ in $\bar{\mathbb{F}}$.
Note that with these conventions we may intersect $\mathbb{F}_d$ and $\mathbb{F}_{d'}$ to get that
$\mathbb{F}_{\gcd(d,d')}$ is their largest common subf\/ield.
We now def\/ine representations of $(Q,\mathbf{d})$ by assigning an $\mathbb{F}_{d_i}$-vector space to each vertex $i\in
I$ and an $\mathbb{F}_{\gcd(d_i,d_j)}$-linear map to each arrow from vertex~$i$ to vertex~$j$.
The f\/inite-dimensional representations of $(Q,\mathbf{d})$ form a~f\/initary, hereditary, Abelian category
$\operatorname{rep}_\mathbb{F}(Q,\mathbf{d})$ where kernels and cokernels are taken vertex-wise, we refer the reader  %, where
to~\cite{rupel1} for more details.

Recall that we work under the assumption that the quiver~$Q$ contains no oriented cycles.
Then the simple representations of $(Q,\mathbf{d})$ may be labeled as $S_i$ ($i\in I$) where the only non-zero vector  %, where
space is $\mathbb{F}_{d_i}$ at vertex~$i$.
The Grothendieck group $\mathcal{K}(Q,\mathbf{d})$ of $\operatorname{rep}_\mathbb{F}(Q,\mathbf{d})$ has
a~$\mathbb{Z}$-basis given by the classes $\alpha_i=[S_i]$ ($i\in I$) of the irreducible representations $S_i$, in this
way we identify the root lattice $\mathcal{Q}$ with $\mathcal{K}(Q,\mathbf{d})$.
Write $[V]$ for the isomorphism class of a~representation $V\in\operatorname{rep}_\mathbb{F}(Q,\mathbf{d})$ and write
$|V|$ for its \emph{dimension vector}, i.e.~the class of~$V$ in $\mathcal{K}(Q,\mathbf{d})$.
Def\/ine the \emph{height} $\operatorname{ht}(V)$ of a~representation~$V$ as the length of any composition series for~$V$.

We def\/ine the \emph{Euler--Ringel form}
$\langle\cdot,\cdot\rangle:
\mathcal{K}(Q,\mathbf{d})\times\mathcal{K}(Q,\mathbf{d})\to\mathbb{Z}$ on generators~by
\begin{gather*}
\langle\alpha_i,\alpha_j\rangle=
\begin{cases}
d_i & \text{if $i=j$,}
\\
-[d_ib_{ij}]_+ & \text{if $i\ne j$.}
\end{cases}
\end{gather*}
The identif\/ication of $\mathcal{Q}$ and $\mathcal{K}(Q,\mathbf{d})$ allows to transfer the bilinear form $(\cdot,\cdot)$
to $\mathcal{K}(Q,\mathbf{d})$.
It then follows from the def\/initions that $(\cdot,\cdot)$ is the symmetrization of the Euler--Ringel form, %--
i.e.~$(\alpha_i,\alpha_j)=\langle\alpha_i,\alpha_j\rangle+\langle\alpha_j,\alpha_i\rangle$ for all $i,j\in I$.
Recall that we write $\alpha_i^\vee=\alpha_i/d_i$.
For a~dimension vector $\mathbf{v}\in\mathcal{K}(Q,\mathbf{d})$ def\/ine ${}^*\mathbf{v}\in\mathcal{K}(Q,\mathbf{d})$~by
\begin{gather*}
{}^*\mathbf{v}=\sum\limits_{i\in I}\langle\alpha_i^\vee,\mathbf{v}\rangle\alpha_i.
\end{gather*}
\setcounter{section}{5}
\setcounter{subsection}{0}
\setcounter{subsubsection}{0}
\setcounter{Theorem}{2}
\setcounter{equation}{0}
% Original source lines 539--585
\subsection{Generalized Feigin homomorphisms}
\label{subsec:hall_gen_feigin}

Def\/ine the \emph{Hall--Ringel bialgebra} $\mathcal{H}(Q,\mathbf{d})=\bigoplus\mathbb{C}[V]$ to be the free
$\mathcal{K}(Q,\mathbf{d})$-graded $\mathbb{C}$-vector space spanned by the isomorphism classes of representations of
$(Q,\mathbf{d})$ with grading $|[V]|=|V|$ given by dimension vector.
The structure constants of the multiplication and comultiplication count certain extensions between representations.
More precisely, for representations~$U$,~$V$, and~$W$, we def\/ine the \emph{Hall number} as the cardinality of the
following set:
\begin{gather*}
\mathcal{F}_{UW}^V=\{R\subset V: R\cong W \ \text{and} \  V/R\cong U\}.
\end{gather*}
For convenience we abbreviate $v=\sqrt{|\mathbb{F}|}$.
\begin{Theorem}[\cite{ringel2}]
%\label{th:hall_mult}
The map $\mu: \mathcal{H}(Q,\mathbf{d})\otimes\mathcal{H}(Q,\mathbf{d})\to\mathcal{H}(Q,\mathbf{d})$ $($written
multiplicatively$)$ given on basis elements~by
\begin{gather*}
[U][W]=\sum\limits_{[V]}v^{\langle |U|,|W|\rangle} |\mathcal{F}_{UW}^V|\cdot[V]
\end{gather*}
defines an associative multiplication on $\mathcal{H}(Q,\mathbf{d})$.
\end{Theorem}
The tensor square $\mathcal{H}(Q,\mathbf{d})\otimes\mathcal{H}(Q,\mathbf{d})$ will be considered as an algebra via the
twisted multiplication def\/ined~by
\begin{gather*}
([U]\otimes[W])([U']\otimes[W'])=v^{(|W|,|U'|)}[U][U']\otimes[W][W'].
\end{gather*}

\begin{Theorem}[\cite{green}]\label{th:hall_comult}
The map $\Delta: \mathcal{H}(Q,\mathbf{d})\to\mathcal{H}(Q,\mathbf{d})\otimes\mathcal{H}(Q,\mathbf{d})$ given on basis
elements~by
\begin{gather*}
\Delta([V])=\sum\limits_{[U],[W]}v^{\langle |U|,|W|\rangle}
\frac{|\operatorname{Aut}(U)||\operatorname{Aut}(W)|}{|\operatorname{Aut}(V)|}\cdot\big|\mathcal{F}_{UW}^V\big|\cdot[U]\otimes[W]
\end{gather*}
defines a~coassociative comultiplication on $\mathcal{H}(Q,\mathbf{d})$ which is compatible with the multiplication on
$\mathcal{H}(Q,\mathbf{d})$ and the twisted multiplication on
$\mathcal{H}(Q,\mathbf{d})\otimes\mathcal{H}(Q,\mathbf{d})$.
\end{Theorem}

The graded \emph{dual Hall--Ringel bialgebra} $\mathcal{H}^*(Q,\mathbf{d})$ has a~basis of delta functions
$\delta_{[V]}$ with grading $|\delta_{[V]}|=-|V|$.
It will be convenient to consider a~slight rescaling $[V]^*$ of the standard dual basis given~by
\begin{gather*}
[V]^*=v^{-\frac{1}{2}\langle V,V\rangle+\frac{1}{2}\sum\limits_{i=1}^nd_iv_i}\delta_{[V]}.
\end{gather*}

\setcounter{section}{5}
\setcounter{subsection}{1}
\setcounter{subsubsection}{0}
\setcounter{Theorem}{9}
\setcounter{equation}{1}
% Original source lines 704--981
\section{Quantum shuf\/f\/le algebras}\label{sec:shuffle}

The quantized coordinate ring $\mathcal{A}_v[N]$ admits an embedding into a~quantum shuf\/f\/le algebra associated to the
$I\times I$ symmetrizable Cartan matrix~$A$ with symmetrizers $\mathbf{d}$.
In this section we introduce our choice of quantum shuf\/f\/le algebra and prove some basic properties.
Better proofs of many of these results can be found in~\cite{leclerc}.
We refer the reader also to \cite{Rosso} for a more general discussion of quantum shuf\/f\/le algebras.


Let $\mathbf{W}=\bigsqcup\limits_{m\ge0} I^m$ denote the set of all words in the alphabet~$I$.
For a~word $\mathbf{i}=(i_1,\ldots,i_m)$ and a~sequence of nonnegative integers $\mathbf{c}\in\mathbb{Z}_{\ge0}^m$ we
will make use of the following simplifying notation:
\begin{gather*}
\mathbf{i}^\mathbf{c}=(i_1^{c_1},\ldots,i_m^{c_m})=(\underbrace{i_1,\ldots,
i_1}_{c_1},\ldots,\underbrace{i_m,\ldots,i_m}_{c_m}),
\end{gather*}
where $i_k$ appears $c_k$ times ($1\le k\le m$).
For $\mathbf{j}=(j_1,\ldots,j_r)\in\mathbf{W}$ we will write $|\mathbf{j}|=\alpha_{j_1}+\dots+\alpha_{j_r}$ for the
\emph{degree} of the word $\mathbf{j}$.
Denote by $\mathbf{W}^\nu$ the set of all words of degree $\nu\in\mathcal{Q}$.

Def\/ine the quantum shuf\/f\/le algebra $g^*$ as the $\mathbb{Z}[v^{\pm\frac{1}{2}}]$-algebra with basis $\mathbf{W}$ and
multiplication given~by
\begin{gather}
\label{eq:shuffle_product}
(j_1,\ldots,j_r)\circ(j_{r+1},\ldots,j_{r+s})=\sum\limits_{\sigma\in\Sigma_{r,s}}v^{\zeta(\sigma)}\mathbf{j}_\sigma,
\end{gather}
where $\mathbf{j}_\sigma=(j_{\sigma_1},\ldots,j_{\sigma_{r+s}})$ and $\zeta(\sigma)=\zeta_{r,s}(\sigma)$ measures both
inversions and non-inversions occurring in~$\sigma$:
\begin{gather*}
\zeta(\sigma)=\frac{1}{2}\sum\limits_{\substack{1\le k\le r
\\
r+1\le \ell\le r+s
\\
\sigma^{-1}_\ell<\sigma^{-1}_k}}(\alpha_{j_k},\alpha_{j_\ell})-\frac{1}{2}\sum\limits_{\substack{1\le k\le r
\\
r+1\le \ell\le r+s
\\
\sigma^{-1}_k<\sigma^{-1}_\ell}}(\alpha_{j_k},\alpha_{j_\ell}).
\end{gather*}

\begin{Remark}
Kleshchev--Ram~\cite{kleshchev-ram1} and Leclerc~\cite{leclerc} use opposite multiplications for their quantum shuf\/f\/le
algebras measuring only inversions and only non-inversions respectively.
The conventions of~\cite{kleshchev-ram1} are claimed to be more convenient for the relationship with the representation
theory of quiver Hecke algebras.

We hope that our choice of multiplication on $g^*$ will prove equally useful in this regard since this is our preferred
quantum shuf\/f\/le multiplication which will f\/igure prominently in the results presented in this note.
\end{Remark}

The following proposition can be obtained from the results of~\cite{leclerc} by rescaling the basis, however for
completeness we provide a~proof.

\begin{Proposition}%\label{prop:shuffle_associativity}
The shuffle product~\eqref{eq:shuffle_product} defines an associative algebra structure on $g^*$.
\end{Proposition}

\begin{proof}
It will be convenient to introduce the subgroup $\Sigma_{r,s,t}$ of $\Sigma_{r+s+t}$ consisting of all shuf\/f\/les of the
sequences $(1,\ldots,r)$, $(r+1,\ldots,r+s)$, and $(r+s+1,\ldots,r+s+t)$, that is
\begin{gather*}
\Sigma_{r,s,t}=\big\{\sigma\in\Sigma_{r+s+t}: \sigma^{-1}_1<\dots<\sigma^{-1}_r, \sigma^{-1}_{r+1}<\dots<\sigma^{-1}_{r+s},
\sigma^{-1}_{r+s+1}<\dots<\sigma^{-1}_{r+s+t}\big\}.
\end{gather*}
There are natural embeddings $\Sigma_{r,s}\hookrightarrow\Sigma_{r,s,t}$ written $\sigma\mapsto\widetilde\sigma$ and
$\Sigma_{s,t}\hookrightarrow\Sigma_{r,s,t}$ written $\sigma\mapsto\widetilde\sigma'$, where
\begin{gather*}
\widetilde\sigma_k=\sigma_k
\quad
\text{for}
\quad
1\le k\le r+s
\qquad
\text{and}
\qquad
\widetilde\sigma_{r+s+k}=r+s+k
\quad
\text{for}
\quad
1\le k\le t,
\\
\widetilde\sigma'_k=k
\quad
\text{for}
\quad
1\le k\le r
\qquad
\text{and}
\qquad
\widetilde\sigma'_{r+k}=r+\sigma_k
\quad
\text{for}
\quad
1\le k\le s+t.
\end{gather*}
Notice that every element of $\Sigma_{r,s,t}$ has a~unique presentation as $\widetilde\sigma\cdot\tau$ with
$\sigma\in\Sigma_{r,s}$ and $\tau\in\Sigma_{r+s,t}$ and a~unique presentation as $\widetilde\sigma'\cdot\rho$ with
$\sigma\in\Sigma_{s,t}$ and $\rho\in\Sigma_{r,s+t}$.
For $\sigma\in\Sigma_{r,s,t}$ def\/ine
\begin{gather*}
\widetilde\zeta(\sigma)=\frac{1}{2}\sum\limits_{\substack{1\le k\le r
\\
r+1\le\ell\le r+s
\\
\sigma_\ell^{-1}<\sigma_k^{-1}}}(\alpha_{j_k},\alpha_{j_\ell}) +\frac{1}{2}\sum\limits_{\substack{1\le k\le r
\\
r+s+1\le\ell\le r+s+t
\\
\sigma_\ell^{-1}<\sigma_k^{-1}}}(\alpha_{j_k},\alpha_{j_\ell}) +\frac{1}{2}\sum\limits_{\substack{r+1\le k\le r+s
\\
r+s+1\le\ell\le r+s+t
\\
\sigma_\ell^{-1}<\sigma_k^{-1}}}(\alpha_{j_k},\alpha_{j_\ell})
\\
\phantom{\widetilde\zeta(\sigma)=}{}
 -\frac{1}{2}\sum\limits_{\substack{1\le k\le r
\\
r+1\le\ell\le r+s
\\
\sigma_k^{-1}<\sigma_\ell^{-1}}}(\alpha_{j_k},\alpha_{j_\ell}) -\frac{1}{2}\sum\limits_{\substack{1\le k\le r
\\
r+s+1\le\ell\le r+s+t
\\
\sigma_k^{-1}<\sigma_\ell^{-1}}}(\alpha_{j_k},\alpha_{j_\ell}) -\frac{1}{2}\sum\limits_{\substack{r+1\le k\le r+s
\\
r+s+1\le\ell\le r+s+t
\\
\sigma_k^{-1}<\sigma_\ell^{-1}}}(\alpha_{j_k},\alpha_{j_\ell}).
\end{gather*}
Then one immediately checks that:
\begin{gather*}
\widetilde\zeta(\widetilde\sigma\cdot\tau)=\zeta(\sigma)+\zeta(\tau)
\qquad
\text{for any $\sigma\in\Sigma_{r,s}$ and $\tau\in\Sigma_{r+s,t}$,}
\\
\widetilde\zeta(\widetilde\sigma'\cdot\rho)=\zeta(\sigma)+\zeta(\rho)
\qquad
\text{for any $\sigma\in\Sigma_{s,t}$ and $\rho\in\Sigma_{r,s+t}$.}
\end{gather*}
We are now ready to prove the proposition.
Indeed, on the one hand we have
\begin{gather*}
\big((j_1,\ldots,j_r)\circ(j_{r+1},\ldots,j_{r+s})\big)\circ(j_{r+s+1},\ldots,j_{r+s+t})
\\
\qquad =\sum\limits_{\sigma\in\Sigma_{r,s}}v^{\zeta(\sigma)}\mathbf{j}_{\sigma}\circ(j_{r+s+1},\ldots,j_{r+s+t})
 =\sum\limits_{\substack{\sigma\in\Sigma_{r,s}\\ \tau\in\Sigma_{r+s,t}}}v^{\zeta(\sigma)+\zeta(\tau)}\mathbf{j}_{\widetilde\sigma\cdot\tau}
\end{gather*}
and on the other hand we have
\begin{gather*}
(j_1,\ldots,j_r)\circ\big((j_{r+1},\ldots,j_{r+s})\circ(j_{r+s+1},\ldots,j_{r+s+t})\big)
\\
\qquad =\sum\limits_{\sigma\in\Sigma_{s,t}}v^{\zeta(\sigma)}(j_1,\ldots,j_r)\circ\mathbf{j}_\sigma
 =\sum\limits_{\substack{\sigma\in\Sigma_{s,t}\\ \rho\in\Sigma_{r,s+t}}}v^{\zeta(\sigma)+\zeta(\rho)}\mathbf{j}_{\widetilde\sigma'\cdot\rho}.
\end{gather*}
But both of these expressions are equal to
$\sum\limits_{\sigma\in\Sigma_{r,s,t}}v^{\widetilde\zeta(\sigma)}\mathbf{j}_\sigma$.
\end{proof}

We will also need a~comultiplication on $g^*$.
This is def\/ined on a~word $\mathbf{h}\in\mathbf{W}$~by
\begin{gather*}
\Delta(\mathbf{h})=\sum\limits_{\mathbf{j}_1,\mathbf{j}_2\in\mathbf{W}:(\mathbf{j}_1,\mathbf{j}_2)=\mathbf{h}}
v^{\frac{1}{2}(|\mathbf{j}_1|,|\mathbf{j}_2|)}\mathbf{j}_1\otimes\mathbf{j}_2.
\end{gather*}
In the following result we consider $g^*\otimes g^*$ with the twisted multiplication
\begin{gather*}
(\mathbf{j}_{1,1}\otimes\mathbf{j}_{1,2})\circ(\mathbf{j}_{2,1}\otimes\mathbf{j}_{2,2})
=v^{(|\mathbf{j}_{1,2}|,|\mathbf{j}_{2,1}|)}(\mathbf{j}_{1,1}\circ\mathbf{j}_{2,1})\otimes(\mathbf{j}_{1,2}\circ\mathbf{j}_{2,2})
\end{gather*}
for all words $\mathbf{j}_{1,1}$, $\mathbf{j}_{1,2}$, $\mathbf{j}_{2,1}$, $\mathbf{j}_{2,2}\in\mathbf{W}$. %$, $
\begin{Lemma}
The map $\Delta:g^*\to g^*\otimes g^*$ defines a~coassociative comultiplication on $g^*$.
Moreover, $g^*$ is a~bialgebra with counit $\epsilon:g^*\to\mathbb{C}$ given by $\varepsilon(\mathbf{j})=0$ whenever
$\mathbf{j}\ne()$ is not the empty word and $\varepsilon\big(()\big)=1$.
\end{Lemma}

\begin{proof}
We will check explicitly the coassociativity and the compatibility of multiplication and comultiplication on $g^*$, the
other axioms of a~bialgebra structure on $g^*$ are straightforward and left to the reader.

We start with the coassociativity of~$\Delta$.
For a~word $\mathbf{h}\in\mathbf{W}$ we have
\begin{gather*}
(\Delta\otimes\operatorname{Id})\Delta(\mathbf{h})
=\sum\limits_{\mathbf{j}_1,\mathbf{j}_2\in\mathbf{W}:(\mathbf{j}_1,\mathbf{j}_2)=\mathbf{h}}
v^{\frac{1}{2}(|\mathbf{j}_1|,|\mathbf{j}_2|)}\Delta(\mathbf{j}_1)
\otimes\mathbf{j}_2
\\
\phantom{(\Delta\otimes\operatorname{Id})\Delta(\mathbf{h})}
=\sum\limits_{\substack{\mathbf{j}_1,\mathbf{j}_2\in\mathbf{W}:\\(\mathbf{j}_1,\mathbf{j}_2)=\mathbf{h}}}
v^{\frac{1}{2}(|\mathbf{j}_1|,|\mathbf{j}_2|)}
\sum\limits_{\substack{\mathbf{j}_{1,1},\mathbf{j}_{1,2}\in\mathbf{W}:\\(\mathbf{j}_{1,1},\mathbf{j}_{1,2})=\mathbf{j}_1}}
v^{\frac{1}{2}(|\mathbf{j}_{1,1}|,|\mathbf{j}_{1,2}|)}(\mathbf{j}_{1,1}\otimes\mathbf{j}_{1,2})
\otimes\mathbf{j}_2
\\
\phantom{(\Delta\otimes\operatorname{Id})\Delta(\mathbf{h})}
=\sum\limits_{\mathbf{j}_{1,1},\mathbf{j}_{1,2},\mathbf{j}_2\in\mathbf{W}:(\mathbf{j}_{1,1},\mathbf{j}_{1,2},\mathbf{j}_2)=\mathbf{h}}
v^{\frac{1}{2}(|\mathbf{j}_{1,1}|+|\mathbf{j}_{1,2}|,|\mathbf{j}_2|)}
v^{\frac{1}{2}(|\mathbf{j}_{1,1}|,|\mathbf{j}_{1,2}|)}\mathbf{j}_{1,1}\otimes\mathbf{j}_{1,2}\otimes\mathbf{j}_2,
\\
(\operatorname{Id}\otimes\Delta)\Delta(\mathbf{h})
=\sum\limits_{\mathbf{j}_1,\mathbf{j}_2\in\mathbf{W}:(\mathbf{j}_1,\mathbf{j}_2)=\mathbf{h}}
v^{\frac{1}{2}(|\mathbf{j}_1|,|\mathbf{j}_2|)}\mathbf{j}_1\otimes\Delta(\mathbf{j}_2)
\\
\phantom{(\operatorname{Id}\otimes\Delta)\Delta(\mathbf{h})}
=\sum\limits_{\substack{\mathbf{j}_1,\mathbf{j}_2\in\mathbf{W}:\\(\mathbf{j}_1,\mathbf{j}_2)=\mathbf{h}}}
v^{\frac{1}{2}(|\mathbf{j}_1|,|\mathbf{j}_2|)}
\sum\limits_{\substack{\mathbf{j}_{2,1},\mathbf{j}_{2,2}\in\mathbf{W}:\\(\mathbf{j}_{2,1},\mathbf{j}_{2,2})=\mathbf{j}_2}}
v^{\frac{1}{2}(|\mathbf{j}_{2,1}|,|\mathbf{j}_{2,2}|)}\mathbf{j}_1\otimes(\mathbf{j}_{2,1}\otimes\mathbf{j}_{2,2})
\\
\phantom{(\operatorname{Id}\otimes\Delta)\Delta(\mathbf{h})}
=\sum\limits_{\mathbf{j}_1,\mathbf{j}_{2,1},\mathbf{j}_{2,2},\in\mathbf{W}:(\mathbf{j}_1,\mathbf{j}_{2,1},\mathbf{j}_{2,2})=\mathbf{h}}
v^{\frac{1}{2}(|\mathbf{j}_1|,|\mathbf{j}_{2,1}|+|\mathbf{j}_{2,2}|)}v^{\frac{1}{2}(|\mathbf{j}_{2,1}|,|\mathbf{j}_{2,2}|)}
\mathbf{j}_1\otimes\mathbf{j}_{2,1}\otimes\mathbf{j}_{2,2}.
\end{gather*}
But both of these are equal to
\begin{gather*}
\sum\limits_{\mathbf{j}_1,\mathbf{j}_2,\mathbf{j}_3,\in\mathbf{W}:(\mathbf{j}_1,\mathbf{j}_2,\mathbf{j}_3)=\mathbf{h}}
v^{\frac{1}{2}(|\mathbf{j}_1|,|\mathbf{j}_2|)}v^{\frac{1}{2}(|\mathbf{j}_1|,|\mathbf{j}_3|)}
v^{\frac{1}{2}(|\mathbf{j}_2|,|\mathbf{j}_3|)}\mathbf{j}_1\otimes\mathbf{j}_2\otimes\mathbf{j}_3,
\end{gather*}
this establishes coassociativity.

To see the compatibility of the shuf\/f\/le multiplication on $g^*$ and~$\Delta$ we consider two words
$\mathbf{j}_1,\mathbf{j}_2\in\mathbf{W}$.
On one hand we have
\begin{gather*}
\Delta(\mathbf{j}_1\circ\mathbf{j}_2)=\sum\limits_{\sigma\in\Sigma_{r,s}}v^{\zeta(\sigma)}\Delta(\mathbf{j}_\sigma)
=\sum\limits_{\sigma\in\Sigma_{r,s}}v^{\zeta(\sigma)}
\sum\limits_{\mathbf{j}_{\sigma,1},\mathbf{j}_{\sigma,2}\in\mathbf{W}:(\mathbf{j}_{\sigma,1},\mathbf{j}_{\sigma,2})=\mathbf{j}_\sigma}
v^{\frac{1}{2}(|\mathbf{j}_{\sigma,1}|,|\mathbf{j}_{\sigma,2}|)}\mathbf{j}_{\sigma,1}\otimes\mathbf{j}_{\sigma,2}.
\end{gather*}
On the other hand we have
\begin{gather*}
\Delta(\mathbf{j}_1)\circ\Delta(\mathbf{j}_2)
=\left(\sum\limits_{\substack{\mathbf{j}_{1,1},\mathbf{j}_{1,2}\in\mathbf{W}:\\(\mathbf{j}_{1,1},\mathbf{j}_{1,2})=\mathbf{j}_1}}
v^{\frac{1}{2}(|\mathbf{j}_{1,1}|,|\mathbf{j}_{1,2}|)}\mathbf{j}_{1,1}\otimes\mathbf{j}_{1,2}\right)\!
\circ\!\left(\sum\limits_{\substack{\mathbf{j}_{2,1},\mathbf{j}_{2,2}\in\mathbf{W}:\\(\mathbf{j}_{2,1},\mathbf{j}_{2,2})=\mathbf{j}_2}}
v^{\frac{1}{2}(|\mathbf{j}_{2,1}|,|\mathbf{j}_{2,2}|)}\mathbf{j}_{2,1}\otimes\mathbf{j}_{2,2}\right)
\\
\qquad
=\sum\limits_{\substack{\mathbf{j}_{1,1},\mathbf{j}_{1,2},\mathbf{j}_{2,1},\mathbf{j}_{2,2}\in\mathbf{W}:\\
(\mathbf{j}_{1,1},\mathbf{j}_{1,2})=\mathbf{j}_1,(\mathbf{j}_{2,1},\mathbf{j}_{2,2})=\mathbf{j}_2}}
v^{\frac{1}{2}(|\mathbf{j}_{1,1}|,|\mathbf{j}_{1,2}|)}v^{\frac{1}{2}(|\mathbf{j}_{2,1}|,|\mathbf{j}_{2,2}|)}
v^{(|\mathbf{j}_{1,2}|,|\mathbf{j}_{2,1}|)}(\mathbf{j}_{1,1}\circ\mathbf{j}_{2,1})\otimes(\mathbf{j}_{1,2}\circ\mathbf{j}_{2,2})
\\
\qquad
=\sum\limits_{\substack{\mathbf{j}_{1,1},\mathbf{j}_{1,2},\mathbf{j}_{2,1},\mathbf{j}_{2,2}\in\mathbf{W}:\\
(\mathbf{j}_{1,1},\mathbf{j}_{1,2})=\mathbf{j}_1,(\mathbf{j}_{2,1},\mathbf{j}_{2,2})=\mathbf{j}_2}}
\sum\limits_{\sigma_1\in\Sigma_{r_1,s_1}}
\sum\limits_{\sigma_2\in\Sigma_{r_2,s_2}}v^{\frac{1}{2}(|\mathbf{j}_{1,1}|,|\mathbf{j}_{1,2}|)}
v^{\frac{1}{2}(|\mathbf{j}_{2,1}|,|\mathbf{j}_{2,2}|)}v^{(|\mathbf{j}_{1,2}|,|\mathbf{j}_{2,1}|)}
\\
\qquad\phantom{=}  \times
v^{\zeta(\sigma_1)}v^{\zeta(\sigma_2)}\mathbf{j}_{\sigma_1,1}\otimes\mathbf{j}_{\sigma_2,2},
\end{gather*}
where we wrote $r_i$ and $s_i$ ($i=1,2$) for the lengths of $\mathbf{j}_{i,1}$ and $\mathbf{j}_{i,2}$ respectively.
There is an obvious bijection between pairs consisting of a~shuf\/f\/le $\sigma\in\Sigma_{r,s}$ together with a~partition
$(\mathbf{j}_{\sigma,1},\mathbf{j}_{\sigma,2})$ of~$\mathbf{j}_\sigma$ and quadruples consisting of partitions of~$\mathbf{j}_1$ and~$\mathbf{j}_2$ together with shuf\/f\/les $\sigma_1\in\Sigma_{r_1,s_1}$ and
$\sigma_2\in\Sigma_{r_2,s_2}$.
Indeed, reading of\/f the contributions in $\mathbf{j}_{\sigma,1}$ and $\mathbf{j}_{\sigma,2}$ coming from $\mathbf{j}_1$
and $\mathbf{j}_2$ produce the desired partitions while the relative positions of these contributions produce the
desired shuf\/f\/les in $\Sigma_{r_1,s_1}$ and $\Sigma_{r_2,s_2}$.

Notice that under this bijection we have
\begin{gather*}
\zeta(\sigma)=\zeta(\sigma_1)+\zeta(\sigma_2)+\frac{1}{2}(|\mathbf{j}_{1,2}|,|\mathbf{j}_{2,1}|)-\frac{1}{2}(|\mathbf{j}_{1,1}|,|\mathbf{j}_{2,2}|)
\end{gather*}
and
\begin{gather*}
(|\mathbf{j}_{\sigma,1}|,|\mathbf{j}_{\sigma,2}|)
=(|\mathbf{j}_{1,1}|,|\mathbf{j}_{1,2}|)+(|\mathbf{j}_{1,1}|,|\mathbf{j}_{2,2}|)
+(|\mathbf{j}_{2,1}|,|\mathbf{j}_{1,2}|)+(|\mathbf{j}_{2,1}|,|\mathbf{j}_{2,2}|).
\end{gather*}
An easy inspection now shows that $\Delta(\mathbf{j}_1\circ\mathbf{j}_2)=\Delta(\mathbf{j}_1)\circ\Delta(\mathbf{j}_2)$
as claimed.
\end{proof}
\setcounter{section}{6}
\setcounter{subsection}{1}
\setcounter{subsubsection}{0}
\setcounter{Theorem}{12}
\setcounter{equation}{2}
% Original source lines 1296--1609
\section{Quantum shuf\/f\/le characters}\label{sec:shuffle_character}

The results of Sections~\ref{sec:feigin},~\ref{subsec:shuffle_gen_feigin}, and~\ref{subsec:hall_gen_feigin} have
established the commutativity of the front two faces in the tetrahedron below:
\begin{gather}\label{dia:feigin_tetrahedron}
\begin{split}&
\includegraphics{Rupel-Fig1}
\end{split}
\end{gather}
%\begin{equation}%
%\label{dia:feigin_tetrahedron}
% \begin{tikzpicture}[description/.style={fill=white,inner sep=2pt}]
%     \matrix (m) [matrix of math nodes, row sep=3em,column sep=2.5em, text height=1.5ex, text depth=0.25ex]
%      { & \mathcal{H}^*(Q,\mathbf{d}) & & \\ \mathcal{A}_v[N] & & & P_\mathbf{i} \\ & & g^* & \\};
%     \path[->,font=\scriptsize]
%      (m-1-2) edge[dashed] node[near start, above=4pt] {$\Omega$} (m-3-3)
%      (m-2-1) edge (m-1-2) edge (m-3-3) edge[line width=6pt,draw=white] (m-2-4) edge node[near start,above] {$\Psi_\mathbf{i}$} (m-2-4)
%      (m-1-2) edge node[above] {$\widetilde\Psi_\mathbf{i}$} (m-2-4)
%      (m-3-3) edge node[below=3pt] {$\overline\Psi_\mathbf{i}$} (m-2-4);
%\end{tikzpicture}
%\end{equation}
Our goal of this section is to complete this diagram, that is def\/ine the \emph{quantum shuffle character} homomorphism
$\Omega: \mathcal{H}^*(Q,\mathbf{d})\to g^*$ making the remaining two faces of the tetrahedron commute.

For a~word $\mathbf{j}=(j_1,\ldots,j_r)$ write $\mathcal{F}_\mathbf{j}(V)$ for the set of f\/lags in~$V$ of type~$\mathbf{j}$:
\begin{gather*}
\mathcal{F}_\mathbf{j}(V)=\{0=V_r\subset V_{r-1}\subset\dots\subset V_1\subset V_0=V: V_{k-1}/V_k\cong S_{j_k}\;\text{for $1\le k\le r$}\}. %\;
\end{gather*}
Then for a~basis vector $[V]^*\in\mathcal{H}^*(Q,\mathbf{d})$ we def\/ine $\Omega([V]^*)\in g^*$ as a~certain generating
function for counting f\/lags in~$V$:
\begin{gather}
\label{eq:Omega}
\Omega([V]^*)=\sum\limits_{\mathbf{j}\in\mathbf{W}} v^{-\sum\limits_{k<\ell}\langle\alpha_{j_\ell},\alpha_{j_k}\rangle}
|\mathcal{F}_\mathbf{j}(V)|\cdot\mathbf{j}.
\end{gather}
This map $\Omega:\mathcal{H}^*(Q,\mathbf{d})\to g^*$ turns out to preserve a~fair amount of structure.

\begin{Theorem}
\label{th:shuffle_character}
The map $\Omega: \mathcal{H}^*(Q,\mathbf{d})\to g^*$ given by~\eqref{eq:Omega} is a~homomorphism of twisted bialgebras.
\end{Theorem}

\begin{Remark}
In general,~$\Omega$ will not be injective.
For example, every member of the $\mathbb{P}^1$-family of indecomposable representations with dimension vector $(1,1)$
for the 2-Kronecker quiver will have the same image under~$\Omega$.
\end{Remark}

\begin{proof}
We begin by showing that~$\Omega$ is an algebra homomorphism.
Recall the rescaling
\begin{gather*}
[V]^*=v^{-\frac{1}{2}\langle |V|,|V|\rangle+\frac{1}{2}\sum\limits_{i=1}^n d_iv_i}\delta_{[V]}.
\end{gather*}
The product $[U]^*[W]^*$ can then be computed in this same basis as
\begin{gather*}
[U]^*[W]^*=v^{-\frac{1}{2}\langle |U|+|W|,|U|+|W|\rangle+\frac{1}{2}\langle |U|,|W|\rangle+\frac{1}{2}\langle
|W|,|U|\rangle+\frac{1}{2}\sum\limits_{i=1}^n d_i(u_i+w_i)}\delta_{[U]}\delta_{[W]}
\\
\phantom{[U]^*[W]^*}
 =\sum\limits_{[V]}v^{-\frac{1}{2}\langle |V|,|V|\rangle+\frac{1}{2}\langle |U|,|W|\rangle+\frac{1}{2}\langle
|W|,|U|\rangle+\frac{1}{2}\sum\limits_{i=1}^n d_iv_i}v^{\langle|U|,|W|\rangle}
\\
\phantom{[U]^*[W]^*=}{}
\times
\frac{|\operatorname{Aut}(U)||\operatorname{Aut}(W)|}{|\operatorname{Aut}(V)|}\cdot\big|\mathcal{F}^V_{UW}\big|\cdot \delta_{[V]}
\\
%\label{eq:hall_basis_prod}
\phantom{[U]^*[W]^*}
 =\sum\limits_{[V]}v^{\frac{3}{2}\langle |U|,|W|\rangle+\frac{1}{2}\langle
|W|,|U|\rangle}\frac{|\operatorname{Aut}(U)||\operatorname{Aut}(W)|}{|\operatorname{Aut}(V)|}\cdot\big|\mathcal{F}^V_{UW}\big|\cdot [V]^*.
\end{gather*}
Thus the equality $\Omega([U]^*[W]^*)=\Omega([U]^*)\Omega([W]^*)$ is equivalent to the following interesting collection
of identities involving f\/lags.
\begin{Lemma}
\label{lem:flag_identity}
For every word $\mathbf{h}\in\mathbf{W}$ we have
\begin{gather}
\nonumber
\sum\limits_{[V]} v^{\frac{3}{2}\langle |U|,|W|\rangle+\frac{1}{2}\langle|W|,|U|\rangle-\sum\limits_{k<\ell}\langle\alpha_{h_\ell},\alpha_{h_k}\rangle}
\frac{|\operatorname{Aut}(U)||\operatorname{Aut}(W)|}{|\operatorname{Aut}(V)|}\big|\mathcal{F}^V_{UW}\big||\mathcal{F}_\mathbf{h}(V)|
\\
\qquad
 =\sum\limits_{\substack{\mathbf{j}_1,\mathbf{j}_2\in\mathbf{W}\\
\sigma\in\Sigma_{\operatorname{ht}(U),\operatorname{ht}(W)}: \mathbf{j}_\sigma=\mathbf{h}}}v^{-\sum\limits_{1\le
k<\ell\le t}\langle\alpha_{j_\ell},\alpha_{j_k}\rangle-\sum\limits_{t+1\le k<\ell\le
r+s}\langle\alpha_{j_\ell},\alpha_{j_k}\rangle+\zeta(\sigma)}|\mathcal{F}_{\mathbf{j}_1}(U)||\mathcal{F}_{\mathbf{j}_2}(W)|,
\label{eq:flag_identity}
\end{gather}
where we abbreviated $t=\operatorname{ht}(U)$.
\end{Lemma}

\begin{proof}
This essentially follows from the compatibility of multiplication and comultiplication in the Hall--Ringel algebra
$\mathcal{H}(Q,\mathbf{d})$.
Using the def\/initions it is easy to compute the following composition of comultiplication and iterated multiplication:
\begin{gather*}
\Delta([S_{h_1}]\cdots[S_{h_{r+s}}])=\sum\limits_{[V]}v^{\sum\limits_{k<\ell}\langle
\alpha_{h_k},\alpha_{h_\ell}\rangle}|\mathcal{F}_\mathbf{h}(V)|\Delta([V])
\\
%\label{eq:comulit_iter_mult}
\qquad
 =\sum\limits_{[U],[V],[W]}v^{\langle|U|,|W|\rangle+\sum\limits_{k<\ell}\langle\alpha_{h_k},\alpha_{h_\ell}\rangle}
\frac{|\operatorname{Aut}(U)||\operatorname{Aut}(W)|}{|\operatorname{Aut}(V)|}\big|\mathcal{F}^V_{UW}\big||\mathcal{F}_\mathbf{h}(V)|\cdot[U]\otimes[V].
\end{gather*}
Theorem~\ref{th:hall_comult} asserts that this is the same as the following iterated multiplication of
comultiplications:
\begin{gather*}
\Delta([S_{h_1}])\cdots\Delta([S_{h_{r+s}}])=([S_{h_1}]\otimes1+1\otimes[S_{h_1}])\cdots([S_{h_{r+s}}]\otimes1+1\otimes[S_{h_{r+s}}])
\\
%\label{eq:iter_mult_comult}
\qquad
 =\sum\limits_{\substack{[U],[W]:\\ \operatorname{ht}(U)+\operatorname{ht}(W)=r+s\\
\sigma\in\Sigma_{\operatorname{ht}(U),\operatorname{ht}(W)}\\
\mathbf{j}_1,\mathbf{j}_2\in\mathbf{W}: \mathbf{j}_\sigma=\mathbf{h}}}v^E
|\mathcal{F}_{\mathbf{j}_1}(U)||\mathcal{F}_{\mathbf{j}_2}(W)|\cdot [U]\otimes[W],
\end{gather*}
with exponent $E=\sum\limits_{1\le k<\ell\le t}\langle\alpha_{j_k},\alpha_{j_\ell}\rangle
+\sum\limits_{t+1\le k<\ell\le r+s}\langle\alpha_{j_k},\alpha_{j_\ell}\rangle
+\sum\limits_{\substack{1\le k\le t\\t+1\le\ell\le r+s\\\sigma_k^{-1}>\sigma_\ell^{-1}}}(\alpha_{j_k},\alpha_{j_\ell})$,
where we abbreviate $t=\operatorname{ht}(U)$.

A few remarks seem necessary to clarify the second equality above.
To obtain a~f\/ixed $[U]$ and $[W]$ we must take from the product
\begin{gather*}
([S_{h_1}]\otimes1+1\otimes[S_{h_1}])\cdots([S_{h_{r+s}}]\otimes1+1\otimes[S_{h_{r+s}}])
\end{gather*}
exactly $\operatorname{ht}(U)$ factors, recorded by $\mathbf{j}_1$, of the form $[S_{h_k}]\otimes1$ and
$\operatorname{ht}(W)$ factors, recorded by $\mathbf{j}_2$, of the form $1\otimes[S_{h_k}]$.
With this notation we are looking at the products $[S_{j_1}]\cdots[S_{j_t}]\otimes[S_{j_{t+1}}]\cdots[S_{j_{r+s}}]$, where
again we have used the abbreviation $\operatorname{ht}(U)=t$.
The f\/irst and second sums in~$E$ as well as the f\/lag coef\/f\/icients should now be transparent.
In order to record the contribution of the twisted multiplication on
$\mathcal{H}(Q,\mathbf{d})\otimes\mathcal{H}(Q,\mathbf{d})$ to the coef\/f\/icient of $[U]\otimes[W]$ we need a~shuf\/f\/le
$\sigma\in\Sigma_{\operatorname{ht}(U),\operatorname{ht}(W)}$ and $j_k=h_{\sigma_k^{-1}}$.
Indeed, the condition $\sigma_k^{-1}>\sigma_\ell^{-1}$ of the last sum in~$E$ records a~product of the form
$(1\otimes[S_{j_k}])([S_{j_\ell}]\otimes1)$ and thus produces the stated contribution.

As previously stated, Theorem~\ref{th:hall_comult} allows us, for a~f\/ixed $[U]\otimes[W]$, to extract the equality below
from which the claim will follow:
\begin{gather}
\sum\limits_{[V]}v^{\langle|U|,|W|\rangle+\sum\limits_{k<\ell}\langle\alpha_{h_k},\alpha_{h_\ell}\rangle}
\frac{|\operatorname{Aut}(U)||\operatorname{Aut}(W)|}{|\operatorname{Aut}(V)|}|\mathcal{F}^V_{UW}||\mathcal{F}_\mathbf{h}(V)|
\nonumber
\\
\qquad
=\sum\limits_{\substack{\sigma\in\Sigma_{\operatorname{ht}(U),\operatorname{ht}(W)}\\
\mathbf{j}_1,\mathbf{j}_2\in\mathbf{W}: \mathbf{j}_\sigma=\mathbf{h}}}
v^E|\mathcal{F}_{\mathbf{j}_1}(U)||\mathcal{F}_{\mathbf{j}_2}(W)|.
\label{eq:green_equality}
\end{gather}
To complete the proof we perform a~few simple manipulations with the sums appearing in the exponents
of~\eqref{eq:flag_identity} and~\eqref{eq:green_equality}.
To start we match the left hand sides of~\eqref{eq:flag_identity} and~\eqref{eq:green_equality} by multiplying both
sides of~\eqref{eq:flag_identity} by $v^{\sum\limits_{k<\ell}\langle
\alpha_{h_\ell},\alpha_{h_k}\rangle-\frac{1}{2}\langle |U|,|W|\rangle-\frac{1}{2}\langle |W|,|U|\rangle}$ and both sides
of~\eqref{eq:green_equality} by $v^{-\sum\limits_{k<\ell}\langle \alpha_{h_k},\alpha_{h_\ell}\rangle}$.
Thus it remains to establish the following identity of exponents.
\begin{Lemma}
\label{lem:exp_equality}
For any $\sigma\in\Sigma_{\operatorname{ht}(U),\operatorname{ht}(W)}$ and words $\mathbf{j}_1$, $\mathbf{j}_2$ satisfying
$\mathcal{F}_{\mathbf{j}_1}(U)\ne\varnothing$, $\mathcal{F}_{\mathbf{j}_2}(W)\ne\varnothing$, and
$\mathbf{j}_\sigma=\mathbf{h}$ we have:
\begin{gather}
\nonumber
\sum\limits_{k<\ell}\langle \alpha_{h_\ell},\alpha_{h_k}\rangle
-\!\sum\limits_{1\le k<\ell\le t}\!\!\langle\alpha_{j_\ell},\alpha_{j_k}\rangle
-\!\sum\limits_{t+1\le k<\ell\le r+s}\!\!\langle\alpha_{j_\ell},\alpha_{j_k}\rangle
+\zeta(\sigma)-\frac{1}{2}\langle |U|,|W|\rangle-\frac{1}{2}\langle|W|,|U|\rangle
\nonumber
\\
\qquad
=-\sum\limits_{k<\ell}\langle \alpha_{h_k},\alpha_{h_\ell}\rangle+\sum\limits_{1\le k<\ell\le t}\!\!\langle\alpha_{j_k},\alpha_{j_\ell}\rangle
+\sum\limits_{t+1\le k<\ell\le r+s}\!\!\!\!\langle\alpha_{j_k},\alpha_{j_\ell}\rangle
+\!\!\!\sum\limits_{\substack{1\le k\le t\\t+1\le\ell\le r+s\\ \sigma_k^{-1}>\sigma_\ell^{-1}}}\!\!\!\!(\alpha_{j_k},\alpha_{j_\ell}). %\!\!
\label{eq:exp_equality}
\end{gather}
\end{Lemma}

\begin{proof}
We f\/irst eliminate the dependence on $\mathbf{h}$ using the following identities which are immediate consequences of the
equality $j_k=h_{\sigma_k^{-1}}$:
\begin{gather*}
\sum\limits_{k<\ell}\langle \alpha_{h_\ell},\alpha_{h_k}\rangle
-\sum\limits_{1\le k<\ell\le t}\langle\alpha_{j_\ell},\alpha_{j_k}\rangle
-\sum\limits_{t+1\le k<\ell\le r+s}\langle\alpha_{j_\ell},\alpha_{j_k}\rangle
\\
\qquad
=\sum\limits_{\substack{1\le k\le t\\t+1\le\ell\le r+s\\ \sigma_k^{-1}<\sigma_\ell^{-1}}}
\langle\alpha_{j_\ell},\alpha_{j_k}\rangle
+\sum\limits_{\substack{1\le k\le t\\t+1\le\ell\le r+s\\ \sigma_\ell^{-1}<\sigma_k^{-1}}}
\langle\alpha_{j_k},\alpha_{j_\ell}\rangle,
\\
-\sum\limits_{k<\ell}\langle \alpha_{h_k},\alpha_{h_\ell}\rangle
+\sum\limits_{1\le k<\ell\le t}\langle\alpha_{j_k},\alpha_{j_\ell}\rangle
+\sum\limits_{t+1\le k<\ell\le r+s}\langle\alpha_{j_k},\alpha_{j_\ell}\rangle
\\
\qquad
=-\sum\limits_{\substack{1\le k\le t\\t+1\le\ell\le r+s\\ \sigma_k^{-1}<\sigma_\ell^{-1}}}
\langle\alpha_{j_k},\alpha_{j_\ell}\rangle
-\sum\limits_{\substack{1\le k\le t\\t+1\le\ell\le r+s\\ \sigma_\ell^{-1}<\sigma_k^{-1}}}
\langle\alpha_{j_\ell},\alpha_{j_k}\rangle,
\end{gather*}
this reduces~\eqref{eq:exp_equality} to showing
\begin{gather}
\nonumber
\sum\limits_{\substack{1\le k\le t\\t+1\le\ell\le r+s\\ \sigma_k^{-1}<\sigma_\ell^{-1}}}\langle\alpha_{j_\ell},\alpha_{j_k}\rangle
+\sum\limits_{\substack{1\le k\le t\\t+1\le\ell\le r+s\\ \sigma_\ell^{-1}<\sigma_k^{-1}}}\langle\alpha_{j_k},\alpha_{j_\ell}\rangle
+\zeta(\sigma)-\frac{1}{2}\langle|U|,|W|\rangle-\frac{1}{2}\langle |W|,|U|\rangle
\\
\qquad
=-\sum\limits_{\substack{1\le k\le t\\t+1\le\ell\le r+s\\ \sigma_k^{-1}<\sigma_\ell^{-1}}}\langle\alpha_{j_k},\alpha_{j_\ell}\rangle
-\sum\limits_{\substack{1\le k\le t\\t+1\le\ell\le r+s\\ \sigma_\ell^{-1}<\sigma_k^{-1}}}\langle\alpha_{j_\ell},\alpha_{j_k}\rangle
+\sum\limits_{\substack{1\le k\le t\\t+1\le\ell\le r+s\\ \sigma_k^{-1}>\sigma_\ell^{-1}}}(\alpha_{j_k},\alpha_{j_\ell}).
\label{eq:exp_equality2}
\end{gather}
Using the dichotomy $\sigma_k^{-1}<\sigma_\ell^{-1}$ or $\sigma_\ell^{-1}<\sigma_k^{-1}$ for $1\le k\le t$ and
$t+1\le\ell\le r+s$ the last sum on the right hand side of~\eqref{eq:exp_equality2} becomes:
\begin{gather}
\label{eq:zeta}
\sum\limits_{\substack{1\le k\le t\\t+1\le\ell\le r+s\\ \sigma_k^{-1}>\sigma_\ell^{-1}}}(\alpha_{j_k},\alpha_{j_\ell})
=\zeta(\sigma)+\frac{1}{2}\sum\limits_{\substack{1\le k\le t\\t+1\le\ell\le r+s}}(\alpha_{j_k},\alpha_{j_\ell}).
\end{gather}
Rewriting
$(\alpha_{j_k},\alpha_{j_\ell})=\langle\alpha_{j_k},\alpha_{j_\ell}\rangle+\langle\alpha_{j_\ell},\alpha_{j_k}\rangle$
and applying~\eqref{eq:zeta} in the right hand side of~\eqref{eq:exp_equality2} gives:
\begin{gather*}
-\sum\limits_{\substack{1\le k\le t\\t+1\le\ell\le r+s\\ \sigma_k^{-1}<\sigma_\ell^{-1}}}\langle\alpha_{j_k},\alpha_{j_\ell}\rangle
-\sum\limits_{\substack{1\le k\le t\\t+1\le\ell\le r+s\\ \sigma_\ell^{-1}<\sigma_k^{-1}}}\langle\alpha_{j_\ell},\alpha_{j_k}\rangle
+\zeta(\sigma)+\frac{1}{2}\sum\limits_{\substack{1\le k\le t\\t+1\le\ell\le r+s}}
(\langle\alpha_{j_k},\alpha_{j_\ell}\rangle+\langle\alpha_{j_\ell},\alpha_{j_k}\rangle)
\\
\qquad
=-\frac{1}{2}\sum\limits_{\substack{1\le k\le t\\t+1\le\ell\le r+s}}
(\langle\alpha_{j_k},\alpha_{j_\ell}\rangle+\langle\alpha_{j_\ell},\alpha_{j_k}\rangle)+\zeta(\sigma)
\\
\qquad
\phantom{=}
+\sum\limits_{\substack{1\le k\le t\\t+1\le\ell\le r+s\\ \sigma_k^{-1}<\sigma_\ell^{-1}}}\langle\alpha_{j_\ell},\alpha_{j_k}\rangle
+\sum\limits_{\substack{1\le k\le t\\t+1\le\ell\le r+s\\ \sigma_\ell^{-1}<\sigma_k^{-1}}}\langle\alpha_{j_k},\alpha_{j_\ell}\rangle
\\
\qquad
=-\frac{1}{2}\langle |U|,|W|\rangle-\frac{1}{2}\langle |W|,|U|\rangle+\zeta(\sigma)
+\sum\limits_{\substack{1\le k\le t\\t+1\le\ell\le r+s\\ \sigma_k^{-1}<\sigma_\ell^{-1}}}\langle\alpha_{j_\ell},\alpha_{j_k}\rangle
+\sum\limits_{\substack{1\le k\le t\\t+1\le\ell\le r+s\\ \sigma_\ell^{-1}<\sigma_k^{-1}}}\langle\alpha_{j_k},\alpha_{j_\ell}\rangle
\end{gather*}
as desired, which completes the proof of Lemma~\ref{lem:exp_equality}.
\end{proof}

This completes the proof of Lemma~\ref{lem:flag_identity}.
\end{proof}

With Lemma~\ref{lem:flag_identity} proven, we may conclude that~$\Omega$ def\/ines a~homomorphism of algebras.
It remains to show that $\Delta\Omega([V]^*)=(\Omega\otimes\Omega)(\Delta([V]^*))$.
The coproduct $\Delta([V]^*)$ can be expanded in the same rescaled basis as
\begin{gather*}
\Delta([V]^*)=v^{-\frac{1}{2}\langle |V|,|V|\rangle+\frac{1}{2}\sum\limits_{i=1}^n d_iv_i}\Delta(\delta_{[V]})
\\
\phantom{\Delta([V]^*)}
 =\sum\limits_{[U],[W]}v^{-\frac{1}{2}\langle |U|+|W|,|U|+|W|\rangle+\frac{1}{2}\sum\limits_{i=1}^n
d_i(u_i+w_i)}v^{\langle|U|,|W|\rangle}\big|\mathcal{F}^V_{UW}\big|\cdot \delta_{[U]}\otimes\delta_{[W]}
\\
%\label{eq:hall_basis_coprod}
\phantom{\Delta([V]^*)}
 =\sum\limits_{[V]}v^{\frac{1}{2}\langle |U|,|W|\rangle-\frac{1}{2}\langle |W|,|U|\rangle} \big|\mathcal{F}^V_{UW}\big|\cdot
[U]^*\otimes[W]^*.
\end{gather*}
Then the identity $\Delta\Omega([V]^*)=(\Omega\otimes\Omega)(\Delta([V]^*))$ is equivalent to the following collection
of identities involving f\/lags.
\begin{Lemma}
For all words $\mathbf{j}_1,\mathbf{j}_2\in\mathbf{W}$ we have
\begin{gather*}
 v^{-\sum\limits_{k<\ell}\langle
\alpha_{h_\ell},\alpha_{h_k}\rangle}v^{\frac{1}{2}(|\mathbf{j}_1|,|\mathbf{j}_2|)}|\mathcal{F}_\mathbf{h}(V)|
 =\sum\limits_{[U],[W]}v^{-\sum\limits_{1\le k<\ell\le r}\langle\alpha_{j_\ell},\alpha_{j_k}\rangle}
v^{-\sum\limits_{r+1\le k<\ell\le r+s}\langle\alpha_{j_\ell},\alpha_{j_k}\rangle}
\\
\qquad
\phantom{=}{}
\times
v^{\frac{1}{2}\langle|U|,|W|\rangle-\frac{1}{2}\langle|W|,|U|\rangle}
\big|\mathcal{F}_{U,W}^V\big|\cdot|\mathcal{F}_{\mathbf{j}_1}(U)|\cdot|\mathcal{F}_{\mathbf{j}_2}(W)|,
\end{gather*}
where we used the word $\mathbf{h}=(\mathbf{j}_1,\mathbf{j}_2)$.
\end{Lemma}

\begin{proof}
This essentially follows from the associativity of multiplication in the Hall--Ringel algebra
$\mathcal{H}(Q,\mathbf{d})$.
Indeed, expanding the products
\begin{gather*}
[S_{h_1}]\cdots[S_{h_{r+s}}]=([S_{j_1}]\cdots[S_{j_r}])\cdot([S_{j_{r+1}}]\cdots[S_{j_{r+s}}])
\end{gather*}
gives
$|\mathcal{F}_\mathbf{h}(V)|=|\mathcal{F}_{U,W}^V|\cdot|\mathcal{F}_{\mathbf{j}_1}(U)|\cdot|\mathcal{F}_{\mathbf{j}_2}(W)|$.
To see the equality of the exponents on either side we note that
\begin{gather*}
\frac{1}{2}(\mathbf{j}_1,\mathbf{j}_2)
=\frac{1}{2}\langle|\mathbf{j}_1|,|\mathbf{j}_2|\rangle+\frac{1}{2}\langle|\mathbf{j}_2|,|\mathbf{j}_1|\rangle
=\frac{1}{2}\langle|U|,|W|\rangle+\frac{1}{2}\langle|W|,|U|\rangle.
\end{gather*}
Combining this with the equality
\begin{gather*}
-\sum\limits_{k<\ell}\langle \alpha_{h_\ell},\alpha_{h_k}\rangle=-\sum\limits_{1\le k<\ell\le r}\langle
\alpha_{j_\ell},\alpha_{j_k}\rangle-\sum\limits_{r+1\le k<\ell\le r+s}\langle
\alpha_{j_\ell},\alpha_{j_k}\rangle-\langle|\mathbf{j}_2|,|\mathbf{j}_1|\rangle
\end{gather*}
completes the proof.
\end{proof}

To f\/inish the proof of Theorem~\ref{th:shuffle_character} we note that~$\Omega$ respects the grading and thus, since the
multiplication on both $\mathcal{H}^*(Q,\mathbf{d})\otimes\mathcal{H}^*(Q,\mathbf{d})$ and $g^*\otimes g^*$ are twisted
by the grading, we may conclude that~$\Omega$ is a~homomorphism of twisted bialgebras.
\end{proof}
\begin{thebibliography}{99}
\bibitem[7]{green}
Green J.A., Hall algebras, hereditary algebras and quantum groups,
  \href{http://dx.doi.org/10.1007/BF01241133}{\textit{Invent. Math.}} \textbf{120} (1995), 361--377.

\bibitem[18]{kleshchev-ram1}
Kleshchev A., Ram A., Homogeneous representations of {K}hovanov--{L}auda
  algebras, \href{http://dx.doi.org/10.4171/JEMS/230}{\textit{J.~Eur. Math. Soc.}} \textbf{12} (2010), 1293--1306,
  \href{http://arxiv.org/abs/0809.0557}{arXiv:0809.0557}.

\bibitem[23]{leclerc}
Leclerc B., Dual canonical bases, quantum shuf\/f\/les and {$q$}-characters,
  \href{http://dx.doi.org/10.1007/s00209-003-0609-9}{\textit{Math.~Z.}} \textbf{246} (2004), 691--732, \href{http://arxiv.org/abs/math.QA/0209133}{math.QA/0209133}.

\bibitem[29]{ringel2}
Ringel C.M., Hall algebras and quantum groups, \href{http://dx.doi.org/10.1007/BF01231516}{\textit{Invent. Math.}}
  \textbf{101} (1990), 583--591.

\bibitem[31]{Rosso}
Rosso M., Quantum groups and quantum shuf\/f\/les,  \href{http://dx.doi.org/10.1007/s002220050249}{\textit{Invent. Math.}} \textbf{133}  (1998),  399--416.

\bibitem[34]{rupel1}
Rupel D., On a quantum analog of the {C}aldero--{C}hapoton formula,
  \href{http://dx.doi.org/10.1093/imrn/rnq192}{\textit{Int. Math. Res. Not.}} \textbf{2011} (2011), 3207--3236.

\end{thebibliography}
\end{document}
